% !TEX program = XeLaTeX
\documentclass[final,t]{beamer}

% HITSZ academic poster. Layout follows the Overleaf HITSZ poster template
% https://www.overleaf.com/latex/templates/harbin-institute-of-technology-shenzhen-hitsz-poster-template/pmfzwyhbsnth
% (LPPL 1.3c). Compile with XeLaTeX or LuaLaTeX.

\usepackage[orientation=portrait,size=a0,scale=1.24]{beamerposter}
\usepackage{ctex}
\usepackage{graphicx}
\usepackage{booktabs}
\usepackage{amsmath,amssymb}
\usepackage{tikz}
\usetikzlibrary{calc}
\usepackage{multicol}
\usefonttheme{serif}

\definecolor{hitblue}{cmyk}{1,0.80,0,0}
\definecolor{hitlightblue}{cmyk}{0.35,0.10,0,0}
\definecolor{hitnavy}{RGB}{11,31,75}
\definecolor{hitpanel}{RGB}{244,247,251}

\setbeamercolor{background canvas}{bg=white}
\setbeamercolor{headline}{fg=white,bg=hitnavy}
\setbeamercolor{block title}{fg=white,bg=hitblue}
\setbeamercolor{block body}{fg=black,bg=hitpanel}
\setbeamertemplate{navigation symbols}{}

\newlength{\colwidth}
\setlength{\colwidth}{0.30\textwidth}
\setlength{\columnsep}{0.8cm}

\begin{document}
\begin{frame}[t]
\begin{tikzpicture}[remember picture,overlay]
  \fill[hitnavy] (current page.north west) rectangle ($(current page.north east)+(0,-11.2cm)$);
\end{tikzpicture}

\vspace*{-0.4cm}
\begin{center}
  {\color{white}\fontsize{56}{64}\selectfont\bfseries
  面向科学代理模型的残差校准方法}\\[0.55cm]
  {\color{white}\fontsize{32}{40}\selectfont
  Residual-Score Calibration of Scientific Surrogate Models}\\[0.7cm]
  {\color{white}\fontsize{24}{30}\selectfont
  陈远\,\textsuperscript{1} \quad 王雪\,\textsuperscript{1} \quad 李明\,\textsuperscript{2}}\\[0.35cm]
  {\color{white!88}\fontsize{20}{26}\selectfont
  \textsuperscript{1}哈尔滨工业大学（深圳）计算机科学与技术学院 \quad
  \textsuperscript{2}哈尔滨工业大学计算学部}\\[0.2cm]
  {\color{hitlightblue}\fontsize{18}{22}\selectfont
  联系：yuan.chen@stu.hit.edu.cn \quad|\quad 中国计算机大会 2026 海报}
\end{center}

\vspace{1.5cm}

\begin{columns}[t]
\begin{column}{\colwidth}
\begin{block}{研究动机}
高保真 PDE 求解器在反问题与不确定性量化中代价高昂。数据驱动代理模型可以缩短计算时间，但残差分布常随工况漂移，导致区间覆盖率明显低于名义水平。

\vspace{0.35em}
本海报给出一种\textbf{事后、单调、可复现}的残差校准流程：不改动冻结代理 $f_\theta$，只在较小的标注校准集上估计运输映射。
\end{block}

\begin{block}{方法}
对校准样本 $\{(x_i,y_i)\}_{i=1}^{n}$ 定义分数
\[
  s_i=\|y_i-f_\theta(x_i)\|_2.
\]
记 $\hat F$ 为经验分布。查询点 $x$ 处 $1-\alpha$ 对称区间为
\[
  f_\theta(x)\pm \hat F^{-1}(1-\alpha).
\]
若校准分数与测试分数可交换，则覆盖率至少为 $1-\alpha$，并带有 $1/(n+1)$ 的有限样本修正。当分数膨胀超过阈值时回退到数值求解器。
\end{block}
\end{column}

\begin{column}{\colwidth}
\begin{block}{实验设置}
\begin{itemize}
  \setlength{\itemsep}{6pt}
  \item Burgers 方程、Darcy 流、反应--扩散
  \item 冻结 Fourier Neural Operator 作为 $f_\theta$
  \item 校准集 $n=200$，目标覆盖率 $90\%$
  \item 对照：未校准残差、分位回归
\end{itemize}
\end{block}

\begin{block}{主要结果}
\renewcommand{\arraystretch}{1.25}
{\small
\begin{tabular}{@{}lccc@{}}
\toprule
问题 & 校准前 & 本方法 & 宽度比 \\
\midrule
Burgers & 0.71 & 0.91 & 1.31 \\
Darcy & 0.64 & 0.89 & 1.55 \\
反应--扩散 & 0.77 & 0.92 & 1.29 \\
\bottomrule
\end{tabular}}

\vspace{0.45em}
覆盖率回到名义水平；宽度增加可控。当代理错过分岔时，分数膨胀检测触发求解器回退。
\end{block}
\end{column}

\begin{column}{\colwidth}
\begin{block}{结论}
残差校准是科学代理模型的低成本保险：不训练新权重，只要求带标注的校准集。该流程已在三类 PDE 上验证，并作为哈工深学位论文与预印本的共享基线。
\end{block}

\begin{block}{致谢与资料}
感谢国家自然科学基金与哈尔滨工业大学（深圳）的支持。\\[0.3em]
模板结构参考 Overleaf \emph{HITSZ Poster Template}（LPPL 1.3c）。\\[0.3em]
{\footnotesize 预印本 / 代码二维码可替换为实际链接。}
\end{block}
\end{column}
\end{columns}

\end{frame}
\end{document}
